% ECONOMICS_PROBLEM: {"id": "M54-529", "title": "Long-run effects of remote and hybrid work", "jel_code": "M54", "source_id": "OP-0027"}
% !TeX program = xelatex
\documentclass[11pt,a4paper]{article}
\usepackage[margin=23mm]{geometry}
\usepackage{fontspec}
\setmainfont{Latin Modern Roman}
\usepackage{amsmath,amssymb,mathtools}
\usepackage{xcolor,enumitem,booktabs,longtable,array,xurl,hyperref}
\definecolor{muted}{HTML}{53636C}
\hypersetup{colorlinks=true,urlcolor=blue,linkcolor=blue}
\setlength{\parindent}{0pt}
\setlength{\parskip}{5pt}
\newcommand{\R}{\mathbb R}
\newcommand{\E}{\mathbb E}
\newcommand{\N}{\mathbb N}
\newcommand{\ind}{\mathbf 1}
\newcommand{\Var}{\operatorname{Var}}
\newcommand{\Cov}{\operatorname{Cov}}
\newcommand{\supp}{\operatorname{supp}}
\DeclareMathOperator*{\argmin}{arg\,min}
\DeclareMathOperator*{\argmax}{arg\,max}
\newcommand{\entrymeta}[3]{{\sffamily\small\color{muted}\textbf{\texttt{#1}}\quad|\quad #2\par #3\par}}
\newcommand{\Problemref}[1]{\texttt{#1}}
\newcommand{\JELref}[2]{\texttt{#2} (JEL \texttt{#1})}
\begin{document}
\section*{Long-run effects of remote and hybrid work}
\textbf{Problem ID:} \texttt{M54-529}\quad \textbf{JEL:} \texttt{M54}\par
% BEGIN ECONOMICS STATEMENT
\label{problem:OP-0027}
\label{atlas:prob:L7}
\noindent\textbf{Original atlas ID:} L7 (entry 27). \textbf{Classification:} Editorial long-horizon policy, interference, and equilibrium problem.

\noindent\textbf{Original atlas question:} How does remote work affect productivity, promotion, geography, sorting, innovation and family formation?

\subsection*{Definitions and assumptions}
Fix a baseline worker cohort, teams, a horizon $H$, and arrangements $r\in\{o,h,f\}$ denoting fully onsite, specified hybrid, and fully remote work. A treatment protocol includes hours, equipment, management, eligibility, and office access. Let $\mathbf r$ be the vector of team assignments and $e_j(\mathbf r)$ an announced exposure mapping for team $j$. Potential outcomes $Q_{jt}(\mathbf r)$ measure team output on a prespecified quality-adjusted scale; $Y_{it}(\mathbf r)$, $P_{it}(\mathbf r)$, and $F_{it}(\mathbf r)$ measure earnings, promotion, and family outcomes for baseline workers, including leavers. Do not assume no interference unless justified; communication and mentoring naturally connect workers.
\subsection*{Formal research question}
For two specified allocation policies $\pi_1,\pi_0$ over assignment vectors, identify
\[
 \tau_H^V=\mathbb E_{\mathbf r\sim\pi_1}
       \left[\sum_{t=0}^{H}\beta^t V_t(\mathbf r)\right]
       -\mathbb E_{\mathbf r\sim\pi_0}
       \left[\sum_{t=0}^{H}\beta^t V_t(\mathbf r)\right],
\]
where $\beta\in(0,1]$ is fixed and $V$ is separately chosen as output, worker earnings, promotion probability, innovation, or another defined outcome. Characterize heterogeneity by pretreatment type and spillover exposure. For economy-wide adoption, construct a separate model $M$ of firm entry, worker sorting, residence, commuting, and office/housing prices, and bound its equilibrium outcome difference under the corresponding policies.
\subsection*{Interpretation and scope}
The finite-cohort experimental estimand and the economy-wide equilibrium effect are different objects. Measuring workers only while they remain in their original firm excludes a potentially central retention and sorting pathway. Innovation proxies require a definition and validation; counts of messages or patents are not automatically productivity. Randomizing teams rather than individuals can identify some spillovers, but only exposure combinations with design support are identified. Families and residential moves may respond over longer horizons than available records, so their effects require explicitly stated follow-up and missing-data assumptions. A statistically insignificant productivity contrast does not demonstrate equality: prespecified equivalence margins and uncertainty intervals are needed for that conclusion.
\subsection*{Source audit and status}
[S1] is a firm-specific experiment; [S2] and [S3] document broader work-from-home patterns. Added [S4] provides a 2024 randomized hybrid-work trial, so short-run performance and retention effects are already studied. The editorial frontier is transport, organizational spillovers, and long-horizon equilibrium outcomes. The atlas's tractability label is not a mathematical claim, and no universal open status is established.

\subsection*{References}
\begin{enumerate}[label=\textnormal{[S\arabic*]},leftmargin=*]
\item Nicholas Bloom, James Liang, John Roberts, and Zhichun Jenny Ying (2015). \emph{Does Working from Home Work? Evidence from a Chinese Experiment}. The Quarterly Journal of Economics 130(1), 165--218. \url{https://academic.oup.com/qje/article-abstract/130/1/165/2337855}. \textit{Locator:} Primary publisher, working-paper, or author record: bibliographic information and abstract; full-text theorem verification is not claimed. \textit{Role:} Randomized work-from-home evidence in a particular firm and occupation.
\item Alexander Bick, Adam Blandin, and Karel Mertens (2023). \emph{Work from Home before and after the COVID-19 Outbreak}. American Economic Journal: Macroeconomics 15(4). \url{https://doi.org/10.1257/mac.20210061}. \textit{Locator:} Primary publisher, working-paper, or author record: bibliographic information and abstract; full-text theorem verification is not claimed. \textit{Role:} Measurement of work-from-home patterns; descriptive evidence.
\item José María Barrero, Nicholas Bloom, and Steven J. Davis (2023). \emph{The Evolution of Work from Home}. Journal of Economic Perspectives 37(4), 23--49. \url{https://doi.org/10.1257/jep.37.4.23}. \textit{Locator:} Primary publisher, working-paper, or author record: bibliographic information and abstract; full-text theorem verification is not claimed. \textit{Role:} Review and measurement of remote-work arrangements.
\item Nicholas Bloom, Ruobing Han, and James Liang (2024). \emph{Hybrid Working from Home Improves Retention without Damaging Performance}. Nature 630, 920--925. \url{https://doi.org/10.1038/s41586-024-07500-2}. \textit{Locator:} Primary publisher, working-paper, or author record: bibliographic information and abstract; full-text theorem verification is not claimed. \textit{Role:} More recent randomized hybrid-work evidence; limits a blanket claim that hybrid performance is unstudied.
\end{enumerate}

\subsection*{Common conventions: Source mathematical conventions}

Unless an entry states otherwise, agent and firm sets are finite. State and action spaces are measurable, strategies and outcome maps are measurable, probability laws are specified, and expectations exist for the targets under discussion. Infinite-horizon discount factors lie in $(0,1)$ and discounted objectives are finite. Norms, tolerances and complexity input sizes must be specified for computational comparisons. Constants or primitives that appear locally are inputs, not empirical facts asserted by this catalogue.

\textbf{Identification.} A statistical model $m\in\mathcal M$ implies an observable law $P_m$ and target $\theta(m)$. At law $P$, its identified set is
\[
\Theta(P;\mathcal M)=\{\theta(m):m\in\mathcal M,\ P_m=P\}.
\]
Point identification holds when this set is a singleton, or equivalently when $P_{m_1}=P_{m_2}$ implies $\theta(m_1)=\theta(m_2)$ for all compatible models. Sharp bounds mean bounds attained, or approached when the boundary is not attained, by compatible models. Writing this set is a definition, not a solution: the substantive task is to establish informative restrictions and derive or estimate the set. An unrestricted-confounding model may give only trivial bounds.

\textbf{Inference.} Identification, estimability and valid uncertainty quantification are separate questions. Fix $\alpha\in(0,1)$. A confidence procedure $C_n$ over a declared law class $\mathcal P$ has asymptotic uniform coverage at level $1-\alpha$ when
\[
\liminf_{n\to\infty}\inf_{P\in\mathcal P}\Pr_P\{\theta(P)\in C_n\}\geq1-\alpha.
\]
For set-identified targets the coverage event must be defined separately, for example coverage of the full identified set. If an entry uses identification-indexed models instead of uniquely indexed laws, the same coverage convention is applied over those models. Pointwise limits do not establish uniform validity, and asymptotic validity does not guarantee finite-sample precision. Sample sizes, dependence structures and weak-identification sequences are part of each application.

\textbf{Counterfactuals.} $Y_i(a)$ is a potential outcome under a specified intervention, including its timing, implementation and versions. It is not $Y_i$ conditional on voluntarily choosing $a$. A full intervention vector permits interference; shorthand scalar treatment notation does not silently impose no interference. Assignment, selection, attrition, anticipation and measurement must be specified. A claim about an instrument's local effect is not automatically a population policy effect.

\textbf{Economic equilibrium.} A counterfactual model includes best responses, constraints, feasibility, market clearing, state transitions and expectations consistent with the stated information. When several equilibria exist, either a declared selector is an input or the target is set-valued. Comparisons across policy regimes must use the stated selection convention. Reduced-form relationships are not presumed invariant after an equilibrium-changing intervention.

\textbf{Optimization and welfare.} Optimization is over the stated feasible set. If attainment is not established, $\sup$ or $\inf$ and $\varepsilon$-optimal choices replace an asserted maximizer or minimizer. For example, with value $V=\sup_{a\in\mathcal A}W(a)<\infty$, an $\varepsilon$-optimal set is $\{a\in\mathcal A:W(a)\geq V-\varepsilon\}$ for $\varepsilon>0$. Benefits, costs, risk and health or environmental outcomes can be aggregated only using declared units and conversion weights. Interpersonal utility weights, the treatment of fiscal transfers and foreign profits, discounting, and the behavioral welfare criterion are explicit inputs. A comparative-static derivative additionally requires existence and appropriate differentiability of the selected counterfactual.

\textbf{Sources.} Local labels [S1], [S2], and so forth restart in every question. Locators identify the relevant abstract, model, section or result at the level actually checked; a bibliographic or abstract-level verification is not represented as inspection of an entire proof. Working papers are distinguished from later publications and changing versions. Source summaries are paraphrased. All URL checks and editorial formulations refer to the audit date stated above.
% END ECONOMICS STATEMENT
\end{document}
